\documentclass[10pt,a4paper]{article}
\newcommand{\DocCode}{ERP-SDD-001}
\usepackage[a4paper,margin=20mm,headheight=16pt]{geometry}
\usepackage[T1]{fontenc}
\usepackage{lmodern}
\usepackage{microtype}
\usepackage{etoolbox}
\usepackage{atbegshi}
\usepackage[table]{xcolor}
\usepackage{graphicx}
\usepackage{array,booktabs,longtable,tabularx}
\usepackage{enumitem}
\usepackage{amsmath}
\usepackage{fancyhdr}
\usepackage{titlesec}
\usepackage{float}
\usepackage{pdflscape}
\usepackage{listings}
\usepackage{xurl}
\usepackage{hyperref}
\definecolor{navy}{HTML}{15324F}
\definecolor{teal}{HTML}{007F86}
\definecolor{mist}{HTML}{EEF5F7}
\definecolor{slate}{HTML}{46586B}
\definecolor{amber}{HTML}{9A5700}
\hypersetup{colorlinks=true,linkcolor=navy,urlcolor=teal,
  pdfauthor={ICT University ERP project team; AI-assisted draft},
  pdflang={en},bookmarksnumbered=true}
\urlstyle{tt}
\AtBeginShipout{\AtBeginShipoutAddToBox{\special{papersize=\the\paperwidth,\the\paperheight}}}
\setlength{\parindent}{0pt}
\setlength{\parskip}{5pt}
\setlength{\emergencystretch}{3em}
\setcounter{tocdepth}{1}
\renewcommand{\arraystretch}{1.18}
\setlist{nosep,leftmargin=*,topsep=4pt}
\titleformat{\section}{\Large\bfseries\color{navy}}{\thesection}{0.7em}{}
\titleformat{\subsection}{\large\bfseries\color{teal}}{\thesubsection}{0.7em}{}
\titleformat{\subsubsection}{\normalsize\bfseries\color{navy}}{\thesubsubsection}{0.7em}{}
\titlespacing*{\section}{0pt}{12pt}{8pt}
\titlespacing*{\subsection}{0pt}{10pt}{4pt}
\pagestyle{fancy}
\fancyhf{}
\fancyhead[L]{\small\color{slate}ICT UNIVERSITY \quad / \quad ERP}
\fancyhead[R]{\small\color{slate}\DocCode\ \textbar\ v1.0}
\fancyfoot[L]{\footnotesize\color{slate}SEN4121 \quad Summer 2026 \quad Review draft}
\fancyfoot[R]{\small\thepage}
\renewcommand{\headrulewidth}{0.5pt}
\newcolumntype{L}[1]{>{\raggedright\arraybackslash}p{#1}}
\newcolumntype{Y}{>{\raggedright\arraybackslash}X}
\newcommand{\code}[1]{\path{#1}}
\newcommand{\source}[1]{\par{\footnotesize\color{slate}\textbf{Evidence:} \path{#1}}\par}
\newcommand{\statusI}{\textbf{\textcolor{teal}{Implemented}}}
\newcommand{\statusP}{\textbf{\textcolor{amber}{Partial}}}
\newcommand{\statusT}{\textbf{\textcolor{slate}{Target}}}
\newcommand{\callout}[2]{\par\smallskip
  \noindent\fcolorbox{teal}{mist}{\parbox{\dimexpr\linewidth-2\fboxsep-2\fboxrule\relax}{
  \textbf{\color{navy}#1}\par #2}}\par\smallskip}
\newcommand{\cover}[2]{
  \hypersetup{pageanchor=false}
  \begin{titlepage}
  \vspace*{8mm}
  {\color{teal}\rule{\textwidth}{3pt}}\par\vspace{8mm}
  {\Large\bfseries\color{navy}ICT UNIVERSITY}\par
  {\large Faculty of Information and Communication Technologies}\par
  \vspace{18mm}
  {\Huge\bfseries\color{navy}#1\par}
  \vspace{6mm}{\Large\color{teal}Multi-tenant University ERP\par}
  \vspace{4mm}{\large Academic \textbullet\ Marketing \& Finance\par
  Administration \& Human Resources}\par
  \vspace{15mm}
  \begin{tabularx}{\textwidth}{@{}L{40mm}Y@{}}
  \toprule
  Document identifier & \DocCode\\
  Version / status & 1.0 / Evidence-grounded review draft\\
  Baseline & Backend \texttt{6e10c54}; frontend \texttt{6cc6859}; reviewed 13 September 2026\\
  Course & SEN4121 -- Large System Environment, Summer 2026\\
  Instructor & Engr.\ Tanwi Nkiamboh\\
  Scope & React frontend and four-service backend, with deployment and integration boundaries\\
  Prepared with & AI assistance using Copilot SDK in VS Code; team review and approval pending\\
  \bottomrule
  \end{tabularx}
  \vfill
  \callout{Reading guide}{#2}
  {\small This document describes observed code, specified requirements and
  unverified targets separately. It is not a certification of production readiness,
  statutory payroll compliance or successful execution of application tests.}
  \end{titlepage}
}
\newcommand{\frontmatter}[1]{
  \pagenumbering{roman}
  \hypersetup{pageanchor=true}
  \phantomsection
  \section*{Document control}
  \addcontentsline{toc}{section}{Document control}
  \begin{tabularx}{\textwidth}{@{}L{29mm}L{33mm}Y@{}}
  \toprule Version & Date & Change\\\midrule
  1.0 & 13 September 2026 & Initial code-grounded specification with editable PlantUML sources and local PDF build.\\
  \bottomrule
  \end{tabularx}
  \subsection*{Review and approval}
  \begin{tabularx}{\textwidth}{@{}L{44mm}Y L{27mm}@{}}
  \toprule Responsibility & Name / signature & Date\\\midrule
  Product / academic representative & Pending team nomination & Pending\\
  Finance and payroll representative & Pending team nomination & Pending\\
  Technical lead / document reviewer & Pending team nomination & Pending\\
  Course submission approval & Pending group review & Pending\\
  \bottomrule
  \end{tabularx}
  \subsection*{Authority and interpretation}
  The supplied four-page SEN4121 project brief is the source of stakeholder
  obligations. Source code, migrations and tests determine what is implemented.
  Existing operational reports are historical evidence, not rerun results.
  The ISO/IEC/IEEE 29148 reference is used as an organizational alignment for
  requirements; no conformance assessment or reproduction of that standard is claimed.
  \subsection*{Academic integrity disclosure}
  This specification and its diagrams were generated with AI assistance from the
  project brief and repository evidence. Group members must review, correct and
  understand the material, disclose this assistance at the defense, and provide
  their own signed originality declaration and individual contribution statements.
  No names, signatures, approvals, measurements or contributions have been invented.
  \subsection*{Pagination policy}
  #1 Main narrative pages are deliberately bounded; extensive requirements,
  interface inventories and diagrams are appendices. The brief limits the
  combined SRS and technical report to 20 pages excluding appendices and diagrams.
  The package budgets six SRS narrative pages, eight SDD narrative pages and three
  front-matter pages per document: 20 combined non-appendix pages including covers,
  control pages and contents. Appendices are explicitly labeled and excluded from
  that total; this is not two separate 20-page allowances.
  \clearpage\tableofcontents\clearpage
}
\newcommand{\diagram}[4][portrait]{
  \clearpage
  \begingroup
  \def\sheetwidth{297mm}\def\sheetheight{420mm}
  \ifstrequal{#1}{wide}{\def\sheetwidth{420mm}\def\sheetheight{297mm}}{}
  \ifstrequal{#1}{large}{\def\sheetwidth{420mm}\def\sheetheight{594mm}}{}
  \ifstrequal{#1}{large-wide}{\def\sheetwidth{594mm}\def\sheetheight{420mm}}{}
  \newgeometry{layoutwidth=\sheetwidth,layoutheight=\sheetheight,margin=18mm}
  \setlength{\paperwidth}{\sheetwidth}
  \setlength{\paperheight}{\sheetheight}
  \setlength{\headwidth}{\textwidth}
  \subsection{#3}
  \begin{figure}[H]\centering
  \includegraphics[width=\textwidth,height=0.84\textheight,keepaspectratio]{diagrams/rendered/#2.png}
  \caption{#3. #4}\label{fig:#2}
  \end{figure}
  {\footnotesize Editable source: \code{diagrams/#2.puml}. This is a large-format
  diagram sheet; use original size or the companion SVG for full-detail review.
  Solid relationships denote
  implemented structure unless otherwise marked; dashed notes distinguish logical,
  external or target elements.}
  \clearpage
  \endgroup
  \restoregeometry
}
\lstset{basicstyle=\ttfamily\footnotesize,breaklines=true,
  frame=single,rulecolor=\color{teal},backgroundcolor=\color{mist},
  columns=fullflexible,keepspaces=true,showstringspaces=false}

\hypersetup{pdftitle={Software Design Document -- ICT University ERP}}
\begin{document}
\cover{Software Design\\Document}{
Eight-page architecture narrative, followed by detailed design decisions,
21 PlantUML views, frontend/API integration tables and an automatically generated
catalogue of 97 service-local routes and 42 persisted models.}
\frontmatter{Read the SRS for stakeholder requirements and acceptance criteria.
This SDD describes both the React frontend and the backend implementation.}
\pagenumbering{arabic}

\section{Architectural baseline and system decomposition}
\subsection{Scope and viewpoint}
This document describes frontend baseline \texttt{6cc6859} in \code{ERP_System}
and backend baseline \texttt{6e10c54} in \code{ERP_System_backend}. All paths are
backend-relative unless prefixed \code{../ERP_System}. The supplied SEN4121 brief
drives mandatory boundaries; the SRS identifies requirements and implementation gaps.
Figures distinguish code-backed structures from conceptual scaling and recovery targets.

\begin{tabularx}{\textwidth}{@{}L{27mm}Y L{36mm}@{}}
\toprule Boundary & Responsibility & Technology / owned data\\\midrule
Web frontend & Session UX, route shell, student academics/payments, staff self-service & React 19, TypeScript, Vite; no private database.\\
Identity & Users, tenants/campuses, Argon2 login, RS256 signing, refresh families & FastAPI; \code{identity_db}.\\
Academic & Enrollment, curriculum, grades/appeals, attendance, exams/reports & FastAPI + outbox worker; \code{academic_db}.\\
Finance & Fees, invoices, payments/ledger, expenses, campaign ROI and summaries & FastAPI + combined worker; \code{finance_db}.\\
HR & Recruitment, employee records, leave, shifts, payroll, assets & FastAPI; \code{hr_db}.\\
Gateway & Single-origin entry, proxy routes, delegated authentication, IP budgets & Nginx; no business persistence.\\
\bottomrule
\end{tabularx}
\subsection{Structural decisions}
Bounded contexts own data and deployment images. APIs do not perform cross-database
joins. Academic acceptance triggers Finance asynchronously, avoiding a synchronous
availability dependency at enrollment time. HTTP serves interactive reads and commands;
the broker provides durable workflow transfer, not arbitrary service discovery.
Compose DNS names and explicit Nginx upstreams constitute the static service registry.

Four databases share one PostgreSQL 17 process: this is logical ownership separation,
not physical high availability. RabbitMQ and the host are additional shared failure
domains. Workers reuse their service image but are separately runnable processes.
There is no shared runtime business-model package; the \code{contracts} directory
shares schemas rather than Python entities.

\subsection{Trace to requirements}
Implements the architectural shape of DATA-01, INT-01--03 and OPS-01.
Figure~\ref{fig:03-components} shows ownership and broker topology;
Figure~\ref{fig:16-deployment} shows actual containers and production prerequisites.
Complete domain administrative UI, full CRUD, scheduled off-host backups and verified
production rollout are not implied by these architectural components.

\clearpage
\section{Frontend architecture and interaction design}
\subsection{Composition and state}
\code{main.tsx} creates one TanStack Query client, with one query retry and no
refetch-on-window-focus. It wraps \code{App}, \code{AuthProvider} and a React Router
browser router. \code{ProtectedRoute} waits for session restoration and redirects
unauthenticated users to login while preserving their pathname. It does not enforce
module roles; feature pages gate their rendered content.

\begin{tabularx}{\textwidth}{@{}L{22mm}L{36mm}Y@{}}
\toprule Route & Visible experience & Current implementation boundary\\\midrule
\code{/login} & Public login form & Zod email/password validation, inline server error, return navigation.\\
\code{/} & Welcome and role & Not a dynamic metrics dashboard.\\
\code{/academic} & Student academics & Enrollment, offering registration, grades/appeals, risk and PDFs; non-student placeholder.\\
\code{/finance} & Student invoices & Payment intent polling and receipts; non-student placeholder.\\
\code{/people} & Staff workspace & Leave, code-based attendance, payslips and notifications; all non-staff roles receive placeholder.\\
\code{/status} & Authenticated status page & Real gateway/service liveness requests, not dependency readiness.\\
\code{/settings} & Placeholder & No implemented institution/user administration UI.\\
\bottomrule
\end{tabularx}
\subsection{Transport and presentation}
Domain service adapters map snake-case HTTP fields into camel-case TypeScript
models. \code{authFetch} attaches the in-memory bearer token and includes cookies;
on 401 it uses a shared in-flight refresh promise and retries the original request
once. \code{authFetchJson} parses error bodies into \code{ApiError}; generic typing
is a compile-time assertion, not runtime schema validation.
Report adapters fetch binary blobs and trigger downloads with temporary object URLs.

Forms use React Hook Form/Zod for login and enrollment, and controlled state in
payment/leave/check-in forms. Shared Button, FormField, Skeleton, EmptyState and
ErrorState components provide consistent feedback. Tailwind styling, local font
assets and reduced-motion-aware transitions support a responsive dark interface.
The shell includes skip navigation, desktop sidebar and mobile menu.

\subsection{Important limits}
Query keys such as \code{my-invoices} and \code{my-payslips} do not include principal
identity; authentication transitions do not clear the shared cache. This requires
explicit account-switch isolation work, not a claim of already completed client
isolation. Selected download, nested-grade and payment-poll failure states are not
fully surfaced. These are documented gaps, not code changes in this deliverable.
Figure~\ref{fig:04-frontend} and Appendix~\ref{app:frontend} detail the integration.

\clearpage
\section{Identity, authorization and trust boundaries}
\subsection{Session lifecycle}
Identity signs ten-minute RS256 JWTs; resource services verify issuer, audience,
signature and expiry using the public key. Only Identity mounts the private key.
Refresh sessions use a SHA-256 hash of an opaque token, default fourteen-day lifetime,
family ID and revocation timestamp. Rotation inserts a successor rather than
overwriting the old row, enabling sequential token-reuse detection and family revocation.
Login tracks failures and locks after five failures for fifteen minutes.

Frontend startup performs refresh then \code{/auth/me}; the resulting principal lives
in React context. Access tokens are held only in a module variable, not persistent
browser storage. Resource 401 responses trigger refresh/retry and an expiry callback.
There is no proactive timer that forces logout at the exact JWT expiry instant.
Logout revokes the presented refresh token; access tokens remain valid until expiry.
Single-runtime refresh deduplication does not serialize refreshes across browser tabs.

\subsection{Gateway boundary}
Nginx public exceptions include health/readiness and login/refresh/logout.
Protected routes use internal \code{/internal/verify} authentication subrequests;
identity-derived headers replace client-supplied tenant/user headers.
Business services still validate the bearer token independently. Default request
budgets are 20/s with burst 40 and login 5/min with burst 5, per source IP.
JSON APIs use the same browser origin; no separate browser-to-service CORS channel
is required.

\subsection{Persistence boundary}
Per-request dependencies set \code{app.current_institution_id} and
\code{app.is_platform_admin} using transaction-local settings. UUID values are
validated before the PostgreSQL SET statement. Forced RLS applies to scoped business
tables and Identity users/refresh sessions, including service-owned tables.
Identity institution/campus metadata is not RLS-enabled. Bootstrap grants database ownership
to the service role and revokes public connection access. A separate migration owner
and non-owner runtime role are not implemented by the current bootstrap.
Outbox/inbox internal tables intentionally do not follow ordinary tenant RLS.
Identity credential resolution and trusted workers use explicit platform context.

\subsection{Deployment qualifications}
Base Compose sets development mode and disables the Secure refresh-cookie flag.
The production override replaces images/key mounts but does not remove those
inherited values. Host TLS and effective environment inspection are mandatory
release gates. The locally public payment callback also remains gateway-protected.
These configuration facts supersede optimistic comments in older documentation.
AUTH-01--08 and INT-02 map to Figures~\ref{fig:10-session-sequence} and
\ref{fig:19-browser-session}.

\clearpage
\section{Academic domain and asynchronous billing design}
\subsection{Enrollment and curriculum}
An enrollment links a student UUID with program, term and campus context.
Only accepted enrollment status exists; there is no submitted/admissions approval
state machine. Student/program/term uniqueness prevents duplicate records. The HTTP
transaction inserts both enrollment and outbox event, then commits without calling
Finance. Student creation is self-only; staff/admin supplies a student ID.

Curriculum separates Program, Course, Term and CourseOffering.
Prerequisite edges reference courses twice; recursive reachability rejects cycles.
Registration checks enrollment ownership and program/term compatibility. A prerequisite
passes if a published assessment percentage average in any offering is at least 50.
Offering capacity is stored but not used to reject registration. Course/program
update/delete UI and APIs are not supplied by the present create/list surface.

\subsection{Teaching, reporting and integrity}
Assigned instructors and admin/super-admin may manage teaching records.
Attendance stores sessions and student presence, but mark entry does not check
registration membership. Grades use decimal persistence, explicit publication and
student-only published visibility. Appeal decisions allow submitted $\rightarrow$
under-review/accepted/rejected and under-review $\rightarrow$ accepted/rejected.
A changed accepted score writes GradeAuditLog with old/new value and actor.
Ordinary score edits are still possible without that appeal audit path.

Exam creation/rescheduling obtains an institution/term advisory transaction lock
and tests interval overlap for room, instructor and shared registered students.
Back-to-back intervals are permitted. Conflict search is term-scoped, not institution-wide
across all terms. PDF transcripts average published percentages; attendance summaries
use all created sessions. At-risk status is attendance below 75 percent OR the average
of the latest two published assessment percentages below 50; fewer than two grades
does not produce a grade signal.

\subsection{Broker transfer}
The academic worker polls every two seconds, up to twenty unpublished rows,
with \code{FOR UPDATE SKIP LOCKED}. Persistent publish uses RabbitMQ confirms before
setting \code{published_at}. Confirm-before-commit crashes may repeat delivery.
Finance prefetches ten events, inserts an inbox event ID and generates an invoice
from its fee schedule in one local transaction. Duplicate event/enrollment paths
acknowledge without another invoice. Missing fees reject to the DLQ; transient
failures requeue. There is no distributed transaction or exactly-once broker guarantee.
ACAD-01--11 and FIN-02 map to Figures~\ref{fig:07-erd-academic},
\ref{fig:11-enrollment-sequence}, \ref{fig:13-appeal-states} and \ref{fig:20-exam-activity}.

\clearpage
\section{Finance, payments and analytical projections}
\subsection{Payment intent and reconciliation}
Finance owns the invoice amount in whole XAF. A PaymentIntent stores invoice ID,
provider reference, amount, mobile number and pending/succeeded/failed status.
A partial unique index permits at most one pending/succeeded intent per invoice.
The API commits the intent before calling the provider, preserving evidence if
external I/O fails. It does not claim success on provider request acceptance.

\code{PaymentGateway} is a Python Protocol implemented by
\code{MtnMomoGateway} and \code{CamerPayTestDoubleGateway}. The fake is stateless:
numbers ending in 0000 fail; otherwise a six-second elapsed interval produces
success. It works across API/worker processes without shared memory.
The MTN adapter uses a token endpoint and collection request/status endpoints
with a fifteen-second HTTP timeout; live sandbox validation is not established.
The poller queries status and does not resubmit a request never received by MTN.

\subsection{ACID financial effect}
Reconciliation locks the intent, short-circuits terminal state and asks the provider.
Failure commits failed intent only. Success locks the invoice and, if not already
paid, changes status, posts equal cash debit and tuition revenue credit under one
\code{posting_id}, and inserts a receipt unique by invoice ID. These changes commit
together. There is no separate ledger-transaction table. Balance is enforced by the
application helper, not a cross-row database CHECK constraint. Receipt references
are UUID columns with invoice uniqueness, not foreign keys.

The worker runs payment reconciliation every three seconds and hourly summary
catch-up alongside RabbitMQ consumption. Student UI polls active payment state at
1.5 seconds; success invalidates the invoice query. A provider callback is only a
status-query hint and currently cannot traverse the public gateway without bearer auth.
Provider I/O inside the reconciliation row-lock lifetime affects contention.

\subsection{Expenses, campaigns and summaries}
Expense creation posts expense debit and cash credit atomically. Campaign ROI is
attributed paid invoice revenue minus campaign cost, divided by cost; zero cost
returns null with a reason. Attribution is a simple converted-lead/student join,
not an exclusive marketing attribution model; duplicate lead links can inflate totals.
Monthly summaries aggregate ledger timestamps in half-open month intervals and
append numbered versions. Hourly catch-up covers fully elapsed missing months
since first ledger activity. Neither campaign revenue nor pending invoice totals
replace posted ledger revenue. FIN-01--12 map to Figures~\ref{fig:05-domain-classes},
\ref{fig:08-erd-finance} and \ref{fig:12-payment-sequence}.

\clearpage
\section{HR workflows and payroll computation}
\subsection{People and self-service}
Recruitment stores positions and candidate stages. Hiring creates Employee and
marks Candidate hired in one commit, leaving identity linkage null. General stage
update accepts enum values without a recruitment transition graph; the hire path
rejects an already-hired candidate. There is no complete employee CRUD or identity
provisioning API. Staff self-service resolves employee from the verified user's UUID.

AttendanceShift belongs to Employee and stores the current QR token ID. Issuing
a token replaces that ID and gives a ten-minute token using a distinct HS256
secret/audience. Check-in validates expiry/signature/current ID, employee ownership
and shift/employee uniqueness. The UI accepts pasted token text; camera/biometric
integration and actual presence verification are absent.

Leave approval rejects self-approval and final-state redecision; decision and
notification commit together for linked users. Notifications are fetched records,
not real-time push. Performance reviews record actor, period, rating and comments.
Asset movement locks Asset and checks aggregate assignment against stock;
employee-specific return balance and employee-reference validation are not additional
guarantees. Administrative screens remain placeholders.

\subsection{Exact payroll algorithm}
Let $G$ be gross monthly XAF, $C$ the CNPS ceiling, $r_e,r_m$ employee/employer rates,
and $d$ the standard deduction. The implementation computes:
\[
 B=\min(G,C),\qquad E=B r_e,\qquad M=B r_m,\qquad
 T=\max(0,(G-E)(1-d)).
\]
For ascending marginal brackets with lower/upper bounds $L_j,U_j$ and rate $r_j$:
\[
 I=\sum_j \max(0,\min(T,U_j)-L_j)r_j,\qquad
 N=G-\operatorname{round}_{HU}(E)-\operatorname{round}_{HU}(I).
\]
The final open bracket has no fixed upper limit.
CNPS, taxable base and tax are reported using whole-franc HALF\_UP rounding;
tax is computed on the unrounded intermediate taxable value.
Employer CNPS is reported separately, not deducted from employee net.

\subsection{Release gate}
A PayrollRun snapshots active-employee payslips for a unique institution/period.
Only a draft backed by a manually verified schedule can be approved; only approved
runs appear in employee self-service. Rate verification is a super-admin attestation,
not validation against official CNPS/DGI publications. The selected version is explicit;
the code does not automatically select the legally applicable version by run date.
No payroll funds or finance event are sent. HR-01--11 map to
Figures~\ref{fig:09-erd-hr}, \ref{fig:14-payroll-activity}, \ref{fig:15-leave-sequence}
and \ref{fig:21-stock-activity}.

\clearpage
\section{Deployment, configuration and delivery pipeline}
\subsection{Local runtime}
Compose creates ten containers: four web APIs, two workers, frontend, gateway,
PostgreSQL and RabbitMQ. API images use Python 3.13; frontend is built with Node 22
and served by Nginx 1.27. Service APIs listen on 8000 internally, frontend/gateway
on 80, database on 5432 and broker AMQP on 5672. Local gateway publishes 8081;
broker management publishes loopback 15672. Domain APIs/databases are not public
host listeners in the base stack. The test overlay changes development test exposure.

The frontend Dockerfile runs \code{npm ci} and the TypeScript/Vite build, then copies
\code{dist} into Nginx. Static Nginx provides SPA deep-link fallback and thirty-day
immutable caching for selected asset extensions. Vite development proxies \code{/api}
to the gateway on 8081; gateway-root \code{/healthz} is not covered by that Vite proxy.
Keys are read-only mounted; production secrets must not use dev key material.

\subsection{Delivery definitions versus execution evidence}
Backend Jenkins defines contract tests, shared test PostgreSQL/key setup, a
four-service lint/type/test/build matrix and database isolation tests.
Coverage is gated at 70 percent. Frontend Jenkins defines install, lint/typecheck,
coverage, application build and image build; its coverage configuration does not
explicitly enforce the same numeric threshold.
Neither root CI pipeline pushes images. The separate backend-owned promotion
pipeline expects already published images and a paired frontend ref.

Promotion specifies candidate contracts/E2E, human approval, image pull, backup gate,
migrations, deploy, smoke and release recording. Controller validation, credentials,
webhooks, registry publication, exact image provenance and live execution remain
unverified. Calling a tag ``immutable'' does not enforce registry immutability,
and the image template uses colon tag syntax rather than a dedicated digest field.

\subsection{Production configuration gates}
The override intends loopback gateway binding, pinned registry images and external
keys. Environment and port merge semantics must be validated, not inferred from
comments: development mode and insecure-cookie setting remain inherited.
Host Nginx/TLS and domain provisioning are external prerequisites.
Readiness endpoints check database connectivity; current deploy and monitoring
loops primarily use liveness and do not establish worker or broker readiness.
OPS-01--03 and UI-05 map to Figure~\ref{fig:16-deployment} and
Appendix~\ref{app:operations}.

\clearpage
\section{Operational qualities, failure handling and evolution}
\subsection{Observability and capacity}
Application logging uses Python logging, plus Nginx access/error logs with request
duration. This is not a deployed central log-aggregation service. A scheduled script
checks gateway/service liveness, any enrollment DLQ messages, backlog above 100,
disk at 85 percent and memory at 90 percent by default. It has no worker heartbeat,
WAL age or backup freshness monitoring. HTTP readiness cannot certify successful
message transfer.

The historical local load report records 25 VUs for 90 seconds: read p95 882.9 ms
and 0.48 percent errors. The proposed target is five minutes, p95 below 500 ms and
error rate below one percent. The latency target is not met. Shared-host contention
and per-request connection creation with SQLAlchemy NullPool are hypotheses,
not proven causes. Query plans and dedicated-host measurements must precede tuning.

\subsection{Recovery and rollback}
The existing isolated PITR drill uses base backup and archived WAL with an explicit
target time, reporting 21.8 seconds to restore a tiny test table. It is not the full
ERP cluster and does not prove production RPO/RTO. Scheduled pgBackRest, encrypted
off-host repository, retention and full event/session reconciliation remain targets.
Deployment backup checks inspect filesystem freshness; the default WAL-age gate
of fifteen minutes is weaker than the five-minute proposed RPO objective.

Rollback is a manual image selection with schema-compatibility review; it must not
be mistaken for database restoration. Production outages require an approved
recovery point, isolated restore, broker/outbox/inbox reconciliation and session
review before traffic returns. Figure~\ref{fig:17-recovery-activity} is explicitly
a target runbook, not a performed exercise.

\subsection{Evolution and verification}
Stateless APIs and shared key/session persistence permit a conceptual replica
strategy, but connection budgets, lock duration, summary-generation races, refresh
concurrency and static upstream registration need validation first.
Single-host replicas do not remove the host, database or broker failure domains.
Figure~\ref{fig:18-scaling} marks this as conceptual.

Prioritize missing required instructor/admin dashboards, CRUD coverage, effective
production settings, credentialed-or-disclosed payment integration and evidence
collection. Preserve event compatibility and explicit ownership during change.
The SRS traceability matrix covers unit/integration/E2E assets; appended design
decisions explain known trade-offs rather than concealing them.

\clearpage
\appendix
\section{Design decisions, consistency and data rationale}\label{app:decisions}
\subsection{Decision register}
\small
\begin{longtable}{@{}L{16mm}L{37mm}L{107mm}@{}}
\toprule ID & Decision & Rationale, consequence and verification\\\midrule\endfirsthead
\toprule ID & Decision & Rationale, consequence and verification\\\midrule\endhead
\bottomrule\endfoot
DD-01 & Four service databases & Domain ownership meets the brief; one cluster reduces local resource cost but shares failure/capacity. Verify database connection isolation.\\
DD-02 & REST plus enrollment events & Interactive API contracts remain simple; outbox avoids enrollment dependence on Finance availability. UI must represent eventual invoice creation.\\
DD-03 & Transactional outbox/inbox & Local commits plus idempotent consumer prevent ordinary lost/duplicate business effects. Publish-before-mark crash replay is expected.\\
DD-04 & RS256 and rotating cookies & Resource services need only public key; refresh-token hashes reduce raw-token exposure. Stateless access tokens cannot be immediately revoked by refresh logout alone.\\
DD-05 & Forced RLS + app context & Database policies add tenant enforcement to API checks; owner roles require FORCE. Internal platform contexts and non-RLS plumbing are explicit exceptions.\\
DD-06 & Protocol-based provider adapter & Real and deterministic providers share request/status contracts; fake is stateless across processes. No live-sandbox evidence is claimed.\\
DD-07 & Whole-XAF money & Integers avoid binary decimal-money error; Decimal supports payroll rates. Browser numbers must stay within JavaScript safe integer bounds; no numeric boundary policy is yet enforced across all surfaces.\\
DD-08 & Financial snapshots & Receipts, payslips and numbered monthly summaries preserve historical output at creation. These intentional derived records are not a claim of strict normalization of every relation.\\
DD-09 & Local advisory/row locks & Exam term lock, payment row locks and asset lock protect specific invariants. Other workflows cannot inherit concurrency guarantees from these examples.\\
DD-10 & React Query + service adapters & Separates server cache, form state and transport shape. Per-principal cache invalidation and full runtime response validation remain gaps.\\
DD-11 & Manual payroll attestation & Prevents unverified schedules from backing approval; requires human statutory evidence. Boolean attestation does not preserve who verified or authoritative source provenance.\\
DD-12 & Single-host Compose & Practical local and constrained-host demonstration; Kubernetes and full monitoring cluster are not prerequisites. Production isolation/TLS/backup still require configuration.\\
DD-13 & Explicit partial delivery & Placeholder pages communicate absence rather than fake analytics. Required dashboards and lifecycle APIs remain acceptance work, not silently excluded scope.\\
\end{longtable}\normalsize

\subsection{Normalization and denormalization}
Institution, campus and user facts are separated in Identity. Academic separates
program, course and term from offering, enrollment, registration and assessment.
Prerequisite edges resolve the course-to-course relationship without repeated lists.
Finance separates fee configuration from issued invoices and payment attempts.
HR separates recruitment, employee, shift, leave, review, schedule and run.
These structures reduce update anomalies and generally follow dependency separation
appropriate to 3NF-oriented transactional design.

The model intentionally retains repeated \code{institution_id} for tenant policy
evaluation. Refresh sessions copy tenant identity from user at issue time.
Invoice amount snapshots a configured fee; payroll payslips snapshot computed amounts;
financial summaries store versioned aggregates. Outbox payloads and payroll brackets
use JSONB/list-shaped data. Such structures need invariants and provenance rather
than a blanket assertion that every relation is strictly in 3NF. No formal
functional-dependency proof or EXPLAIN comparison is supplied in the baseline.

\subsection{Key constraints and ownership exceptions}
\begin{itemize}
\item \textbf{Identity:} institution slug, user email and refresh-token hash uniqueness;
nullable user institution/campus for platform use.
\item \textbf{Academic:} enrollment student/program/term; course program/code within institution;
prerequisite edge; offering course/term; grade assessment/student; exam one-per-offering.
\item \textbf{Finance:} fee tenant/program/term; invoice tenant/enrollment; inbox event primary
key; provider reference; partial active-intent uniqueness; receipt invoice uniqueness;
summary tenant/period/version.
\item \textbf{HR:} payroll tenant/period; payslip tenant/employee/period; attendance shift/employee.
Some same-database identifiers intentionally lack foreign keys, including payslip employee,
asset-movement employee and receipt references.
\end{itemize}
Cross-service user/student/campus/program/enrollment IDs remain opaque UUIDs.
An ERD dotted relationship is a logical lookup, not a database-enforced FK.
The generated dictionary records only declared column metadata; migrations remain
authoritative for composite/partial indexes and RLS policies.

\subsection{Transaction and concurrency matrix}
\begin{tabularx}{\textwidth}{@{}L{35mm}Y Y@{}}
\toprule Workflow & Atomic unit / lock & Residual qualification\\\midrule
Enrollment & Enrollment + outbox commit & Duplicate pre-check alone is not the concurrency boundary; DB uniqueness is final.\\
Outbox publish & Row locks, confirms, publication timestamp & Confirm and DB commit are separate durability systems; replay possible.\\
Invoice consumption & Inbox + invoice commit & Missing fees roll back and dead-letter; unsupported-version enforcement incomplete.\\
Payment settlement & Intent/invoice row locks; status + ledger + receipt & External status I/O extends lock lifetime.\\
Expense & Expense + balanced ledger commit & Balance is application-checked, not a cross-row database constraint.\\
Exam & Institution/term advisory transaction lock & No cross-term resource check.\\
Grade appeal & Decision + changed score + audit commit & No explicit row lock in decision path; direct grade edits bypass appeal audit.\\
Leave & Decision + notification commit & Concurrent approval race needs independent acceptance testing.\\
Payroll run & Run + all payslips commit & Statutory validation/manual approval are separate operations.\\
Assets & Asset row lock + movement commit & Aggregate stock protected; per-employee over-return is not independently checked.\\
\bottomrule
\end{tabularx}

\clearpage
\section{Detailed frontend integration and UX coverage}\label{app:frontend}
\subsection{Frontend-to-API mapping}
\small
\begin{longtable}{@{}L{31mm}L{64mm}L{65mm}@{}}
\toprule Component / service & HTTP interaction & Refresh behavior and limitation\\\midrule\endfirsthead
\toprule Component / service & HTTP interaction & Refresh behavior and limitation\\\midrule\endhead
\bottomrule\endfoot
AuthProvider & POST \code{/auth/refresh}; GET \code{/auth/me} & Restore on mount; principal set after both succeed.\\
LoginPage & POST \code{/auth/login}; GET \code{/auth/me} & Form validation precedes request; redirects to remembered pathname.\\
AppShell & POST \code{/auth/logout} & Access memory cleared by service finally; shared query cache is not cleared.\\
EnrollmentPage & GET academic programs/terms/enrollments; POST enrollments & Invalidates enrollments; new enrollment invoice polls at 2 s. Non-student queries can start before the placeholder branch.\\
InvoiceCell & GET finance invoices filtered by enrollment ID & Null means invoice not yet generated; only just-created enrollment has interval polling.\\
CoursesPanel & GET courses, offerings, registrations, own at-risk; POST registration & Invalidate registrations and at-risk on successful registration.\\
CourseGrades & GET offering assessments, then assessment grades; POST grade appeal & Lazy queries when expanded; submitted-appeal indication is local component state, not loaded appeal history.\\
Academic reports & GET transcript / attendance-summary & Blob download with fixed file names; rejected downloads lack a complete UI error flow.\\
FinancePage & GET finance invoices & Student-only displayed content; invalidates own invoices on payment success.\\
PayForm & POST invoice payment-intent; GET payment-intent & Local intent ID; 1.5 s poll when data status is pending; failure allows retry. No reload restoration of pending intent selection.\\
Receipt & GET invoice receipt & Offered for paid row; PDF download.\\
LeaveSection & GET leave requests/mine; POST leave requests & Invalidates own leave list; no automatic live approval updates.\\
AttendanceSection & GET attendance/mine; POST attendance/check-in & Manual token entry; invalidate attendance after success.\\
PayslipsSection & GET payroll/payslips/mine & Displays saved whole-XAF gross/CNPS/IRPP/net; no approval UI.\\
NotificationsSection & GET notifications/mine; POST notification/read & Invalidates notification query on read; no timed polling or push.\\
SystemStatusPage & GET gateway and service \code{healthz} paths & Five liveness cards; page reload rechecks. No readiness or worker state.\\
\end{longtable}\normalsize
All business paths above begin with \code{/api/v1}; prefixes are abbreviated for
readability. Identity uses \code{/auth}, HR \code{/hr}, Finance \code{/finance}
and Academic \code{/academic}. See service adapter files for exact path assembly.

\subsection{UI role matrix}
\begin{tabularx}{\textwidth}{@{}L{31mm}L{26mm}L{26mm}Y@{}}
\toprule Page & Student & Staff & Admin / super-admin\\\midrule
Overview / Status & Visible & Visible & Visible; overview is welcome-only.\\
Academic & Functional self-service & Placeholder & Placeholder.\\
Finance & Functional payment portal & Placeholder & Placeholder.\\
People & Placeholder & Functional own workspace & Placeholder despite backend admin capabilities.\\
Settings & Placeholder & Placeholder & Placeholder.\\
\bottomrule
\end{tabularx}
All navigation items are rendered to authenticated users. This is explicit incomplete
feature delivery, not a permission model implemented through hiding links.
The frontend's four role values match backend roles, but department-specific
role separation is not implemented.

\subsection{Response/error contracts and data lifecycle}
HTTP adapters translate fields, not domain rules. Pydantic is server validation;
Zod covers selected input forms. The generic JSON helper uses type assertions and
does not validate arbitrary server responses at runtime. A structured \code{detail}
object can be reduced to an unhelpful string by generic error formatting.
Payment-poll request errors may leave a waiting message; individual grade-query
errors can resemble ``Not graded''. Mark-read mutation errors and download promise
failures need dedicated user feedback.

Query cache uses one application-level client, while access identity changes in
AuthProvider. A robust next revision should cancel/clear principal-owned queries
on logout/session replacement and include stable principal scope where appropriate.
That is a target design, not an implemented property of this baseline.

\subsection{Accessibility and responsive verification scope}
Source supplies visible form labels, focus styles, a skip link, loading status
indicators and reduced-motion support on login/status animations. Manual attendance
code entry avoids dependence on a camera. Existing E2E assets check login, keyboard
flow, invalid credentials and reduced motion, plus a shared authenticated staff
session across overview, people, finance placeholder and status at four viewports:
360x800, 768x1024, 1440x900 and 1920x1080.
The suite deliberately avoids reusing one rotating cookie in many browser contexts.
It is not coverage of every student/admin workflow or full WCAG conformance.
Desktop logout placement, mobile focus/menu behavior and missing error paths need
manual review. No fresh E2E run is claimed by this documentation build.

\clearpage
\section{Message, API and numerical contracts}\label{app:contracts}
\subsection{EnrollmentAccepted.v1 envelope}
\begin{lstlisting}
{
  "event_id": "<outbox UUID>",
  "event_type": "academic.enrollment_accepted",
  "schema_version": 1,
  "tenant_id": "<institution UUID>",
  "occurred_at": "<UTC ISO-8601 timestamp>",
  "correlation_id": "<workflow UUID>",
  "data": {
    "enrollment_id": "<UUID>",
    "student_id": "<UUID>",
    "program_id": "<UUID>",
    "term_id": "<UUID>",
    "campus_id": "<UUID>"
  }
}
\end{lstlisting}
The placeholders are explanatory, not a schema-valid fixture. Use the checked-in
\code{contracts/events/examples/enrollment-accepted.v1.example.json} for executable
contract validation. RabbitMQ routing key carries \code{.v1}; envelope event type
does not. Exchange \code{erp.events} is durable topic; queue
\code{finance.enrollment_accepted} binds \code{academic.enrollment_accepted.*}.
Rejected messages route via direct \code{erp.dlx} to
\code{finance.enrollment_accepted.dlq}.

\textbf{Consumer qualification:} the runtime parses selected fields and UUIDs but
does not fully validate schema version/type, envelope shape and every schema field.
Contract-example tests are separate from runtime consumer validation. Broker
authentication/topology are the trust boundary; the envelope is not cryptographically
signed. Ordering across workers is not a global business guarantee.

\subsection{API examples and status semantics}
\begin{lstlisting}
POST /api/v1/academic/enrollments
Authorization: Bearer <access-token>
Content-Type: application/json
{"program_id":"<UUID>","term_id":"<UUID>","campus_id":"<UUID>"}

POST /api/v1/finance/invoices/<UUID>/payment-intents
{"payer_msisdn":"<synthetic test number>"}

POST /api/v1/hr/leave/requests
{"starts_on":"2026-10-05","ends_on":"2026-10-06","reason":"Example"}
\end{lstlisting}
Academic 201 means enrollment accepted, not invoice paid or generated.
Finance intent 201 means locally persisted attempt, not provider settlement.
HR approval responses record business decisions, not external notification delivery.
404 may deliberately conceal an unauthorized owned resource. The gateway can
generate 429 without the application handler running. Readiness can return 503.
All examples use placeholders or synthetic dates and contain no real credentials.

\subsection{At-risk and finance examples}
At-risk attendance uses present records divided by all sessions of an offering.
Missing marks do not count present. The grade signal averages the latest two
published normalized scores by assessment creation order; one grade is insufficient.
For latest percentages 40 and 55, mean 47.5 flags risk even though both individual
grades are not failures. Scores 40 and 60 average 50 and do not trigger that signal.
Exactly 75 percent attendance is not below threshold.

For campaign cost 100,000 XAF and attributed paid tuition 150,000 XAF,
ROI is 0.5, conventionally displayed as 50 percent. This arithmetic example does
not correct duplicate lead attribution in the query. Monthly net is posted tuition
revenue credits minus expense debits; it is not an accrual balance sheet or a tax return.

\subsection{Illustrative payroll calculation, not statutory rates}
For a deliberately synthetic test schedule with gross 100,000, ceiling 80,000,
employee rate 0.04, employer rate 0.08, deduction 0.30 and a single 10-percent
tax bracket: employee CNPS is 3,200; employer CNPS 6,400; taxable base 67,760;
IRPP 6,776; net 90,024 XAF. These are invented arithmetic-test inputs,
\emph{not published Cameroonian rates}. Real usage requires verified effective-period
CNPS/DGI sources and review of monthly versus annual bracket conventions.

\clearpage
\section{UML and architecture diagram suite}\label{app:diagrams}
This suite exceeds the brief's required use-case, class, two sequence, ERD and
deployment views. All diagrams are editable PlantUML. PNG renders are embedded
for portable LaTeX builds; SVG renders accompany the package for zoomable inspection.
Diagram sheets use A3 or A2 page sizes so dense architecture and sequence views
are not reduced to unreadable A4 thumbnails; narrative and tables remain A4.
Selected-field ERDs must be read with the full generated dictionary and migrations.
SRS and SDD share context/use-case sources to avoid contradictory duplicate artwork.
\diagram{01-context}{System context}{Frontend and backend share a product boundary;
external provider and operator prerequisites are explicitly identified.}
\diagram{02-use-cases}{Actor and use-case model}{Implemented APIs and self-service
screens are distinguished from administrative UI gaps.}
\diagram[wide]{03-components}{Runtime component architecture}{Database ownership, API
routing and the broker-mediated cross-service flow are explicit.}
\diagram{04-frontend}{Frontend composition and transport}{Provider, router, page,
adapter, cache and token responsibilities reflect the React source.}
\diagram[large-wide]{05-domain-classes}{Implementation class collaboration}{Real protocol,
ORM, dataclass and module-function concepts; no invented shared service framework.}
\diagram{06-erd-identity}{Identity entity relationship model}{Nullable tenant/campus
context and refresh-family storage are shown with local foreign keys.}
\diagram{07-erd-academic}{Academic entity relationship model}{Curriculum, enrollment,
teaching records and outbox are distinct; dotted outbox linkage is not a foreign key.}
\diagram{08-erd-finance}{Finance entity relationship model}{Payment and receipt
associations, ledger grouping and analytical projections distinguish logical from FK relationships.}
\diagram[wide]{09-erd-hr}{HR entity relationship model}{Employee linkage, payroll snapshots,
leave notifications and stock movements show local and unenforced references.}
\diagram[large]{10-session-sequence}{Session authentication sequence}{Gateway delegation,
independent validation, rotation and sequential reuse rejection are separate steps.}
\diagram[large]{11-enrollment-sequence}{Enrollment-to-invoice sequence}{Two local ACID
transactions surround at-least-once broker transfer and explicit failure paths.}
\diagram[large]{12-payment-sequence}{Payment settlement sequence}{Durable intent precedes
provider I/O; reconciliation atomically applies the financial effect.}
\diagram{13-appeal-states}{Grade appeal state machine}{Only implemented transitions
are shown; accepted corrections have explicit score and audit guards.}
\diagram[large]{14-payroll-activity}{Payroll generation and approval activity}{Pure calculation,
snapshot creation and manual statutory verification are separate responsibilities.}
\diagram{15-leave-sequence}{Leave decision and notification sequence}{Decision
segregation and in-app notification persistence are shown without fictitious email delivery.}
\diagram[large-wide]{16-deployment}{Container deployment view}{All ten Compose containers,
volumes and network paths are distinguished from unverified production intentions.}
\diagram{17-recovery-activity}{Target full-system recovery activity}{This is a proposed
operational procedure, not the already-run single-table PITR drill.}
\diagram[wide]{18-scaling}{Conceptual scaling design}{Replica groups are design candidates;
shared infrastructure and concurrency limits remain explicit.}
\diagram{19-browser-session}{Browser session state machine}{Startup restoration and
reactive 401 handling are modeled without a nonexistent proactive expiry timer.}
\diagram{20-exam-activity}{Exam scheduling conflict activity}{The actual
institution/term lock and interval conflict checks define the protected invariant.}
\diagram{21-stock-activity}{Aggregate stock movement activity}{The flow protects
aggregate on-hand quantity, not an unimplemented per-employee return constraint.}

\clearpage
\section{Generated service-local route catalogue}\label{app:routes}
Generated by \code{inventory.py} from Python syntax trees without importing services,
connecting to a database or reading secrets. It enumerates actual decorated handlers
and router prefixes. Health and internal routes are service-local, not necessarily
public gateway routes. Return annotations may name ORM classes where no explicit
response model is declared. Access policy is explained in the SRS role matrix and
the owning handler; the table is not a runtime OpenAPI export.
\input{generated/endpoints}

\clearpage
\section{Generated persisted-model dictionary}\label{app:models}
The catalogue lists mapped-column attributes, Python ORM annotations and explicit
column-level metadata. It intentionally excludes ORM relationship properties.
An omitted \code{nullable} keyword does not prove nullability; consult SQLAlchemy
annotation inference and the migration DDL. Composite/partial uniqueness, CHECKs,
indexes and RLS policies must be read from model table arguments and migrations.
The four ERDs provide navigable structure; this appendix provides field detail.
\input{generated/models}

\clearpage
\section{Setup, validation and operational handover}\label{app:operations}
\subsection{Clean local setup sequence}
Prerequisites: Docker with Compose v2 supporting the override tags used by this
repository; sibling frontend checkout; Python 3.13 for service checks; Node 22 for
frontend local builds. Choose an isolated local environment and synthetic records.
The following commands are an operator guide, not commands executed by the document build.
\begin{lstlisting}
# From ERP_System_backend; preserve an existing .env.
Copy-Item .env.example .env
Set-Location services\identity
python -m venv .venv
.\.venv\Scripts\python.exe -m pip install -r requirements-dev.txt
.\.venv\Scripts\python.exe ..\..\scripts\generate_dev_jwt_keys.py
Set-Location ..\..
docker compose up -d --build
docker compose exec identity python -m alembic upgrade head
docker compose exec academic python -m alembic upgrade head
docker compose exec finance python -m alembic upgrade head
docker compose exec hr python -m alembic upgrade head
docker compose exec identity python -m scripts.seed
docker compose exec academic python -m scripts.seed
docker compose exec finance python -m scripts.seed
docker compose exec hr python -m scripts.seed
\end{lstlisting}
Do not overwrite an existing environment file without reviewing it. Use the matching
interpreter for your platform, rather than blindly copying a machine-specific executable.
Seed outputs are development credentials and fixtures, not real accounts.
Verify seed identifier alignment, fee schedules and employee/user linkage before
demonstrating dependent workflows. Never run seed scripts against production records.
Start with gateway \code{http://localhost:8081}; inspect each API's \code{readyz},
not only gateway liveness. Workers have no Compose health check; verify event
transfer by enrolling a synthetic student and observing a Finance invoice.

\subsection{Existing test/build entry points}
\begin{tabularx}{\textwidth}{@{}L{32mm}Y@{}}
\toprule Scope & Existing commands / conditions\\\midrule
Backend service & In each service: \code{ruff check .}, \code{mypy app}, \code{pytest --cov=app --cov-fail-under=70}; install that service's dev requirements and provide migrated PostgreSQL/key configuration as its fixtures require.\\
Contracts & In \code{contracts}: install declared requirements; run \code{pytest}.\\
Isolation & Run \code{tests/integration/test_db_isolation.py} against the isolated configured test cluster.\\
Frontend & \code{npm ci}, \code{npm run lint}, \code{npm run typecheck}, \code{npm run test:coverage}, \code{npm run build}.\\
Frontend E2E & \code{npm run test:e2e}; requires live gateway, dedicated seeded credentials and browser installation. Reusing a rotating cookie across independent contexts is invalid.\\
Documentation & Run \code{docs/specifications/build.ps1} with local Java/PlantUML, Python and Tectonic paths; see package README.\\
\bottomrule
\end{tabularx}
Documentation validation does not mutate runtime data or substitute for these
application tests. Test evidence must identify both frontend and backend revisions.

\subsection{Operational handover gates}
\begin{enumerate}
\item Confirm domain, certificate, effective environment, cookie flags, listener binding
and unique secrets before opening a production origin.
\item Provision an image publication process; record frontend/backend refs and actual
image digests. Validate Jenkins syntax and credentials against the real controller.
\item Establish off-host backup, WAL continuity, retention and monitoring; test a
representative full-cluster restore and event reconciliation.
\item Validate provider mode and payroll rate sources; keep simulations explicit.
\item Resolve mandatory UI/API gaps and collect real coverage, EXPLAIN/load and issue-board evidence.
\item Review legacy comments and historical reports for drift; source and verified
effective runtime configuration remain the authoritative implementation evidence.
\end{enumerate}

\subsection{Documentation verification contract}
The build regenerates route/model tables, renders PlantUML locally to PNG/SVG and
compiles both PDFs with Tectonic. It fails on diagram/compiler errors, undefined
references or overfull boxes. PDF pagination, extractable text and rendered diagram
legibility are checked separately during preparation. No source or diagram text
is submitted to a public PlantUML server; downloaded tools/typesetting packages
do not require sending repository content.

\subsection{References and disclosure}
Primary stakeholder source is the supplied SEN4121 Summer 2026 brief.
Implementation references are backend \texttt{6e10c54}, frontend \texttt{6cc6859},
and the exact source paths in the SRS and generated appendices.
Historical measurements come only from \code{docs/performance.md} and
\code{docs/backup-recovery.md}. The diagrams and narrative are AI-assisted original
documentation of this project, not copied standard text. Team review, signature,
individual contribution statements and disclosure at defense are still required.
\end{document}
